\section{7.2 트리 하나에서 숲으로: 랜덤 포레스트}
\begin{frame}{트리 하나에서 \kb{숲}(랜덤 포레스트)으로}
  \begin{itemize}
    \item 트리 하나는 데이터를 \textbf{완벽히 외워버리기}(과적합) 쉽다
    \item \kb{랜덤 포레스트}(Random Forest, Breiman 2001)의 아이디어:
      서로 조금씩 다른 트리를 \textbf{여러 개} 만들고,
      \textbf{다수결}로 예측한다
  \end{itemize}
  \vskip0.8em
  \begin{block}{핵심 직관}
    \textbf{``완전히 같은 실수를 하는 전문가 100명''}보다
    \textbf{``서로 다른 실수를 하는 전문가 100명''}의 평균이 훨씬 안정적이다.
    두 가지 무작위성이 트리들을 서로 덜 닮게 만들어, 개별 트리의 과적합
    성향이 다수결로 \textbf{상쇄}된다.
  \end{block}
\end{frame}
